% Данный файл распространяется под лицензией CC BY 4.0. Текст лицензии размещён на https://creativecommons.org/licenses/by/4.0/
% (c) Кафедра системного программирования, 2025

\documentclass[a4paper]{article}

\usepackage[a4paper, top=8mm, bottom=8mm, left=8mm, right=8mm]{geometry}

\usepackage{polyglossia}
\setdefaultlanguage[babelshorthands=true]{russian}
\setotherlanguage{english}

\usepackage{fontspec}
\setmainfont{FreeSerif}
\newfontfamily{\russianfonttt}[Scale=0.7]{DejaVuSansMono}

\usepackage[tiny, compact]{titlesec}

\usepackage{titling}
\setlength{\droptitle}{-1cm}
\pretitle{\begin{center}\begin{bfseries}\Large}
\posttitle{\par\end{bfseries}\end{center}}
\preauthor{\begin{center}\normalsize}
\postauthor{\par\end{center}\vspace{-1.8cm}}

\usepackage{hyperref}
\usepackage{bookmark}
\usepackage{csquotes}
\usepackage{upquote}
\usepackage{amsmath}
\usepackage{amssymb}

\title{Разработка интерпретатора языка MiniML на OCaml}

\author{Белимов Максим Андреевич}

\date{}

\begin{document}

\maketitle

\begin{flushright}
    Группа: \emph{26Б71}

    Кафедра: \emph{Кафедра системного программирования}

    Научный руководитель: \emph{Смирнов Кирилл Константинович}
    
    Консультант: \emph{---}
    
    Номер семестра практики: \emph{2}
\end{flushright}

\section{Постановка задачи}

В рамках курсовой работы требовалось реализовать интерпретатор языка miniML — функционального языка семейства ML. Язык поддерживает целочисленные и булевы выражения, \texttt{unit}, кортежи, переменные, лямбда-абстракции (\texttt{\textbackslash x y. \ldots}), связывание имён (\texttt{let}) и рекурсивные определения (\texttt{let rec}), условный оператор с ленивым вариантом (\texttt{if}/\texttt{lif}), обработку исключений (\texttt{try}/\texttt{with}), декларации суммарных типов (\texttt{type}) с сопоставлением по образцу (\texttt{case}), а также арифметические, логические и операции сравнения.

\textbf{Цель работы} — углублённое изучение устройства функциональных языков через практическую реализацию полного интерпретатора на OCaml: от лексического анализа до статической типизации и вычислений.

\section{Описание предлагаемого решения}

Архитектура — конвейер: \textbf{лексер → парсер → AST → тайпчекер → компиляция в λ-исчисление → CBV-редуктор}, обёрнутый REPL-оболочкой.

\subsection{Лексер и парсер}

Лексер реализован как \textbf{детерминированный конечный автомат}: \texttt{parseToken} возвращает пару (длина лексемы, токен) и поддерживает числа, булевы литералы, идентификаторы, ключевые слова, операторы (включая словесные эквиваленты \texttt{and}/\texttt{or}/\texttt{xor}/\texttt{not}) и скобки. Парсер — метод \textbf{рекурсивного спуска} с иерархией приоритетов (let → fun → if → =,!= → || → \^{} → \&\& → сравнения → +,- → *,/ → унарные → аппликация → атомы). Все бинарные операторы правой ассоциативности, аппликация — левой (реализуется итеративным сбором аргументов). Многоместные \texttt{let} и лямбды десахаризуются в одинарные. Полная грамматика приведена в \texttt{docs/GRAMMAR.md}.

\subsection{Основные конструкции}

\texttt{expr} покрывает узлы AST: константы, переменные, \texttt{Let}/\texttt{Let\_rec}, \texttt{Lambd}, \texttt{App}, \texttt{If}, \texttt{Try}, \texttt{Tuple}, бинарные/унарные операции, а также \texttt{Constr} и \texttt{Case} для суммарных типов. Декларация \texttt{type T = C1 of T1 | C2 of T2 in \ldots} вводит конструкторы в области видимости; в \texttt{case} проверяется полнота ветвей на этапе разбора. Конструктор может ссылаться на объявляемый тип (самоссылка), что делает вариантные типы \textbf{рекурсивными} (например, \texttt{type Nat = Zero | Succ of Nat} — натуральные числа) и позволяет кодировать списки, деревья и другие итерируемые структуры данных; такие типы представляются как \texttt{RecT}.

\subsection{Реконструкция типов и тайпчекер}

Тайпчекер реализует \textbf{вывод типов по Хиндли—Милнеру} без аннотаций. Для каждого выражения строится система уравнений на типах (\texttt{set\_equations}), которая затем решается \textbf{унификацией} с проверкой на occurrences. Типы представлены как \texttt{TInt}, \texttt{TBool}, \texttt{TUnit}, \texttt{TArrow}, \texttt{TTuple}, рекурсивные \texttt{RecT} и схемы \texttt{Scheme}. Поддерживаются встроенные примитивы \texttt{nth} (обращение к элементу кортежа) и \texttt{raise}.

\subsection{Let-полиморфизм}

Для \texttt{let x = v in body} тип значения обобщается (\texttt{generalize}): оставшиеся свободные типовые переменные не связываются в окружении, образуя \textbf{схему типа} $\forall\ldots.\,\tau$ (\texttt{Scheme}). При использовании переменной схема \textbf{инстанцируется} (\texttt{instantiate}) свежими типовыми переменными. Это обеспечивает полиморфизм: одна и та же функция (например, \texttt{let id = \textbackslash x. x in \ldots}) может применяться к значениям разных типов.

\subsection{Компиляция в расширенное λ-исчисление}

AST транслируется (\texttt{Lambda.ast\_to\_term}) в λ-термы (\texttt{Fun}, \texttt{App}, \texttt{Var} с \textbf{индексами де Брауна}). Булевы значения кодируются \textbf{по Чёрчу} ($\lambda x.\lambda y.x$ и $\lambda x.\lambda y.y$), \texttt{if} и логические операции выражаются через них, \texttt{let} свёртывается в аппликацию, \texttt{let rec} — через \textbf{Z-комбинатор} неподвижной точки. Нарушение чистоты вносит лишь обработка ошибок: примитив \texttt{raise} представлен константой \texttt{Error}, а \texttt{try}/\texttt{with} — конструктором \texttt{Try}.

\subsection{CBV-редуктор и REPL}

Редуктор (\texttt{Interpreter.eval}) выполняет \textbf{call-by-value}-редукцию: подстановка с избеганием захвата (\texttt{subst}, \texttt{shift}), примитивы сопоставляются по образцу, деление на ноль даёт \texttt{Error}. Значения, тип и результат печатаются в REPL.

\section{Эксперименты}

Реализовано трёхуровневое тестирование: модульные (Alcotest) — токены, лексер, парсер; property-based (QCheck) — roundtrip $\texttt{parse}(\texttt{tokenize}(\texttt{expr\_to\_string}(\mathit{ast}))) = \texttt{Ok}(\mathit{ast})$ на десятках тысяч случайных AST; интеграционные (cram) — сценарии работы REPL. Покрытие кода (bisect\_ppx) около 83\%.

\section{Заключение}

Реализован полный интерпретатор miniML: лексер-ДКА, парсер рекурсивного спуска, тайпчекер с реконструкцией типов и let-полиморфизмом, компиляция в расширенное λ-исчисление с CBV-редукцией. Работа демонстрирует применение конечных автоматов, FIRST-множеств, унификации, подстановки с избеганием захвата и кодирования Чёрча.

\textbf{Ограничения:} строгая CBV-стратегия без редукции под λ, ограниченная диагностика ошибок, отсутствие оптимизации хвостовой рекурсии.

\textbf{Репозиторий:} \url{https://github.com/Maxllon/MiniML}

\begin{thebibliography}{9}
    \bibitem{tapl} Пирс Б. {\itshape Типы в языках программирования}. М.: Лямбда Пресс, Добросвет, 2011. 656+xxiv~с.
    \bibitem{dragon} Aho A.V., Lam M.S., Sethi R., Ullman J.D. {\itshape Compilers: Principles, Techniques, and Tools}. 2nd ed. Addison-Wesley, 2006.
    \bibitem{ifpl} Peyton~Jones S.L. {\itshape The Implementation of Functional Programming Languages}. Prentice Hall, 1987.
    \bibitem{sml97} Milner R., Tofte M., Harper R., MacQueen D. {\itshape The Definition of Standard ML (Revised)}. MIT Press, 1997.
\end{thebibliography}

\end{document}
